\chapter{Future Work: Addressing the Depths of Nested Nonlinearity}

\label{chapter_future_work}

The contributions of this thesis are fundamental in nature and as such are intended to open new directions for future research beyond the solutions to the case studies and example problems presented herein.  The techniques developed and demonstrated apply to a broad class of nonlinear functions and specifically allow detailed and exact analysis of these functions which otherwise elude or impede commensurately detailed analysis with previous state of the art techniques.  While neural networks serve as the primary motivation and the studies of Chapter \ref{chapter_network_replication} demonstrate utility in such applications, new results can be obtained when studying a wide variety of other systems using these techniques, as demonstrated in the studies of Chapter \ref{chapter_additional_case_studies}.  Building off these contributions, we motivate and propose several new studies which are uniquely enabled by this work spanning fields as diverse as general mathematics, neural networks, naval radar system design, and nonlinear circuit design and characterization.  This list of studies provides just a small sample of the space of possible studies which can now be executed employing the mathematics, methodology, and developments detailed herein.

\section{Motivations for Addressing Nested Nonlinearity}
This work addresses directly the fundamental mathematical difficulties of converting a nonlinear function such as $\sech(v)$, $\tanh(v)$, $\sech^2(v)$, or $\relu(v)$,  into a polynomial having exact coefficient formulas, adjustable domain of convergence, monic form, and other properties.  Using prior state of the art methods, this challenge is exacerbated and often intractable when treating nested nonlinear forms such as
\begin{align}
	f(x) = \varphi(a_n \varphi( a_{n-1} \varphi( \cdots a_1 \varphi(a_0 x + b_0) + b_1\cdots ) + b_{n-1}) + b_n). \label{eq_nested_nonlinearity}
\end{align}
We elaborate on the motivations for addressing these forms here and propose studies addressing them in different applications in the following section.  Such forms naturally arise in many applications where chains of nonlinear functions are employed either to achieve a desired result or due to the unavoidable nonlinearities present in a problem which is otherwise conventionally treated via linear analysis.  A first example of the former case includes deep neural networks, where nested layers of basic nonlinear functions allow modeling a more intricate nonlinear function or a decision boundary which is not otherwise captured by a linear model or represented by conventional mathematical forms.

It should be noted that as per the well known universal approximation theorem only one layer of nonlinear activation is strictly required to separate a nonlinear decision boundary or learn a nonlinear function on a closed interval \cite{cybenko_approximation_1989}.  However, the universal approximation theorem provides no guarantees that a single layer of nonlinear functions will yield the most efficient model, or even a practically implementable model of the function to be learned.  As shown by studies in neural network explainability which break open the networks to analyze signals in the latent space (e.g., from Olah et al.), the depth of the network is what enables earlier layers to learn features which are then used to formulate decision mechanisms comprised of meaningful algorithms having recurring motifs similar to those occurring in biological systems \cite{olah2020zoom,olah2017feature}.

The motivation for analyzing nested nonlinearity in neural networks is therefore strong since such nesting is required to achieve state of the art performance in many problems.  While it has been determined through empirical analysis that deep neural networks contain recurring motifs comprising meaningful algorithms and circuits, there has not been a formal analysis conducted to prove when these recurring circuits are equivalent to known functions or to prove when the network has derived some new circuit.  Further, the development of a systematic method for identification of a minimum model, represented as a neural network or a more general form, for some nonlinear process is still an open problem in the study of neural networks and nonlinear systems in general.  One way that such questions could be addressed is to analyze the recurring circuits via known circuit, systems, and mathematical analysis techniques, enabled through the use of the AMITE polynomials.  For example, Olah et al. report the occurrence of a high-low frequency detector which is likely used by an image classifier to find edges.  The input-output relationship of such a circuit could be analyzed rigorously via Fourier methods.  However, application of Fourier expansion or application of a Fourier transform to forms such as Eq. \ref{eq_nested_nonlinearity} without any modified expansion techniques applicable to these forms is an intractable challenge.  For example, even the following Fourier integral, wherein the integrand represents a drastically simplified version of \ref{eq_nested_nonlinearity}, eludes analysis:
\begin{align}
	\int_{0}^{\infty} \! 
	    \tanh(\tanh(v)) \sin(\xi v)
	\, \mathrm{d} v .
\end{align}
Attacking this integral via the methods of Chapter \ref{chapter_additional_case_studies} would require reducing $\tanh(\tanh(v))$ to a polynomial form via methods analogous to an extended version of those detailed in Chapter \ref{chapter_improved_expansions}.  While truncated expansions and well-known error propagation methods such as Taylor models could be used, these methods also come with the well-known challenge of managing the increase in the error interval width as the errors propagate through each nonlinear layer.  The lack of a means for performing analysis by directly expanding the nested functions themselves such that the challenge of error propagation can be avoided is the main fundamental mathematical gap to be addressed which is common to the variety of motivating applications to be discussed.

A further example would include neural networks applied to problems involving signal / time series processing.  In this case, it is likely that the network could learn functions which have some representation as an equation tied to physics, even if that representation is not readily expressible in the language of conventional mathematics using known elementary and special functions.  Consider that in an image recognition problem, the decision boundary between an object of class A and an object of class B is defined by the properties of the objects themselves as learned by the network from the data available.  However, consider in contrast that neural networks used in a time series or signal processing problem could be learning a relationship which is indeed based in the fundamental physics of the problem and has some mathematical form possibly not yet notated in the canon of Bateman, Prudnikov, Gradshteyn and Ryzhik, and other catalogers of such forms.

Consider, for example, a neural network used to predict the behavior of a dynamical system.
Such a network could find application in object tracking wherein it would be advantageous to predict the difference between the estimated state (e.g., via track filtering of measured states) of a dynamical system and the actual state of the dynamical system when the system is exhibiting some self-driving non-constant acceleration.  In a trivial case wherein the state estimation technique employed is a polynomial smoothing filter and the system's non-constant acceleration is sinusoidal, an analytical expression for this residual in terms of a general hypergeometric function can be derived via the methods of Appendix \ref{appendix_hyper_form}.  In a non-trivial case, it is hypothesized that the form could be expressed in terms of more general functions such as the Meijer-G, the Fox-H, the more general still I and J functions, or perhaps some other forms unknown still.  Such a result would have broad reaching mathematical implications, especially if an algorithmic method could be derived for converting a string of nested functions into an equivalent Meijer-G or other representation.   Since the Meijer-G and related general functions are expressed in terms of integrals, success in this endeavor would first allow the analytic integration of a broad class of nested functions which have no references in any known integration tables and also provide a promising direction towards a general algorithm for analytic integration, which has eluded mathematicians and computer scientists for decades.

A further example still which motivates the analysis of nested nonlinear functions includes electrical circuits involving chains of nonlinear processes.  Such an example embodies the case where the nonlinearity of the system may be inherent in the limitations of the components employed and the constraints (e.g., size, weight, power, cost) of the application, rather than an explicitly desired part of the system driven by its design goals.  Consider that a complex signal processing system will often contain several stages of mixing, amplification, filtering, and other operations.  All of these operations are usually designed using conventional linear methodology which assumes that the components are operating in a regime sufficiently linear for such approximations to apply.  However, due to intrinsic nonlinearity and intrinsic variability of dynamic environments, these linear approximations will not always apply in practice. Additionally, systems can dynamically move in and out of linear operating regimes, adding a time varying nature to nonlinear effects.  Consider that components can vary within specified tolerances, vary with temperature, have memory effects and coupling effects, be disturbed by signals injected via outside vectors due to environment and malicious actors, and be subject to further environmental effects on any components which serve as transducers between the system and the outside environment.  Such effects create a problem even more complex than that of the nested nonlinearity of the neural network, since a wide variety of heterogeneous nonlinear functions may be nested, preventing repeated application of the same techniques to untangle them and further motivating the need for increased generality.

\section{Mathematical Approach}

The technical approach for addressing the gaps motivated above rests first on a fundamental mathematical approach for handling nested nonlinear functions and executing analysis which would be trivial for single layers of nonlinear functions (e.g., Fourier analysis, polynomial expansion, etc.) on deeply nested nonlinear functions.  The fundamental mathematical approach which could allow analysis of forms such as Eq. \ref{eq_nested_nonlinearity} rests on three key concepts which we elaborate on here.  First, we require the development of systematic method of iteratively applying the methods of Chapter \ref{chapter_improved_expansions} to untangle layers of nonlinearity.  Second, we require a method for exploiting the known forms of the error functions, also derived in Chapter \ref{chapter_improved_expansions} and elaborated in Appendix \ref{appendix_expansion_error_detail} to prevent the errors from propagating through layers of nested nonlinearity.  A method must be developed for accounting for the errors in a region and absorbing them into the coefficients of the polynomial without sacrificing analytic form.  Finally, we require a method of relating the results obtained back to known forms which can be analyzed such as the hypergeometric functions, Meijer-G, Fox-H, I and J, or other yet to be named forms.  If it is possible to apply these methods iteratively, to account for errors as they propagate nonlinearly, and to extract from the result a standard analytic form which can be compared to other analytic forms, it is hypothesized that a general method for handling nested nonlinear functions could be derived.

First, the proposed method for iteratively applying the expansions is as follows.  The expansion methods of Chapter 3 rely on taking a modified Fourier transform of the function to be expanded, exploiting the duality of the Fourier transform to express the function as a modified Fourier integral, and then expanding the kernel of the transform itself to yield a polynomial expansion.  This creates a challenge for handling forms such as $\tanh(\tanh(v))$ because the integral $\int \! \tanh(\tanh(v))\sin(\xi v) \, \mathrm{d}v$ is intractable.  While the Fourier transform of $\tanh(\tanh(v))$ cannot be used directly, the methods of Chapter \ref{chapter_improved_expansions} can be employed to first transform the outer nonlinear function into a polynomial i.e., $\tanh(v) \simeq \sum_{m=0}^{M} T_{2m+1}(M) v^{2m+1}$ revealing the slightly more tractable form
\begin{align}
	\int_{0}^{\infty} \! 
	\tanh(\tanh(v)) \sin(\xi v)
	\, \mathrm{d} v =
	\lim_{M \to \infty} \sum_{m=0}^{M} T_{2m+1}(M) \int_0^{\rho(M)} \! \tanh^{2m+1}(v) \sin(\xi v) \, \mathrm{d}v.
\end{align}
The resulting form yields a fascinating new pattern to be analyzed.  The first few entries are given as:
\begin{align}
\begin{split}
	&\int \! \tanh(v) \sin(\xi v) \, \mathrm{d}v \\
	&\quad = \frac{1}{4} e^{-i v \xi} \left(-e^{2 v}\right)^{-\frac{i \xi}{2}} \left\{ 
	    i e^{2 i v \xi} \left(\mathrm{B}_{-e^{2 v}}\left(\frac{i \xi}{2}+1,0\right)+\mathrm{B}_{-e^{2 v}}\left(\frac{i \xi}{2},0\right)\right) \right. \\
	&\quad\quad\quad \left. -i \left(-e^{2 v}\right)^{i \xi} \left(\mathrm{B}_{-e^{2 v}}\left(1-\frac{i \xi}{2},0\right)+\mathrm{B}_{-e^{2 v}}\left(-\frac{i \xi}{2},0\right)\right)\right\},
\end{split} \\
\begin{split}
	&\int \! \tanh^3(v) \sin(\xi v) \, \mathrm{d}v \\
	&\quad = \frac{1}{16} e^{-i v \xi} \left\{
	    -2 i \left(\xi^2-2\right) e^{2 i v \xi} \left(-e^{2 v}\right)^{-\frac{i \xi}{2}} \left(\mathrm{B}_{-e^{2 v}}\left(\frac{i \xi}{2}+1,0\right)+\mathrm{B}_{-e^{2 v}}\left(\frac{i \xi}{2},0\right)\right) \right. \\
	    &\quad\quad\quad + 2 i \left(\xi^2-2\right) \left(-e^{2 v}\right)^{\frac{i \xi}{2}} \left(\mathrm{B}_{-e^{2 v}}\left(1-\frac{i \xi}{2},0\right)+\mathrm{B}_{-e^{2 v}}\left(-\frac{i \xi}{2},0\right)\right) \\
	    &\quad\quad\quad  \left. -4 \xi e^{2 i v \xi} \tanh (v)-4 \xi \tanh (v)-4 i e^{2 i v \xi} \text{sech}^2(v)+4 i \text{sech}^2(v) \vphantom{\frac{B}{B}} \right\},
\end{split} \\
\begin{split}
	&\int \! \tanh^5(v) \sin(\xi v) \, \mathrm{d}v \\
	&\quad = \frac{1}{96} e^{-i v \xi} \left\{ \vphantom{\frac{B}{B}} 2 \xi \left(\xi^2-20\right) \left(1+e^{2 i v \xi}\right) \tanh (v)+2 i \left(\xi^2-24\right) \left(-1+e^{2 i v \xi}\right) \text{sech}^2(v)\right. \\
	&\quad\quad\quad + i \left(\xi^4-20 \xi^2+24\right) e^{2 i v \xi} \left(-e^{2 v}\right)^{-\frac{i \xi}{2}} \left(\mathrm{B}_{-e^{2 v}}\left(\frac{i \xi}{2}+1,0\right)+\mathrm{B}_{-e^{2 v}}\left(\frac{i \xi}{2},0\right)\right) \\
	&\quad\quad\quad -i \left(\xi^4-20 \xi^2+24\right) \left(-e^{2 v}\right)^{\frac{i \xi}{2}} \left(\mathrm{B}_{-e^{2 v}}\left(1-\frac{i \xi}{2},0\right)+\mathrm{B}_{-e^{2 v}}\left(-\frac{i \xi}{2},0\right)\right) \\
	&\quad\quad\quad \left. +12 i \left(-1+e^{2 i v \xi}\right) \text{sech}^4(v)+4 \xi \left(1+e^{2 i v \xi}\right) \tanh (v) \text{sech}^2(v) \vphantom{\frac{B}{B}} \right\}.
\end{split}
\end{align}
Obtaining a general form for this pattern would be a useful first step in untangling a doubly nested nonlinear chain of functions.  If this worked out to a tractable form so that the Fourier transform of $\tanh(\tanh(v))$ could easily be computed, this Fourier transform could be used to expand $\tanh(\tanh(v))$ via the methods of Chapter \ref{chapter_improved_expansions}.  It is an open question if the method could be chained to handle functions such as $\tanh(\tanh(\tanh(v)))$.

Next, the method for accounting for the errors proceeds as follows.  We first note from Chapter \ref{chapter_improved_expansions} that the error forms for $\tanh(v)$, $\relu(v)$, and $\sech(v)$ reduce to convenient and accurate closed form approximations valid within the domain of validity $(-V, V)$.  These approximations themselves,
\begin{align}
	I_M^{\{\sech\}}(v) &\simeq
	\frac{
		4 e^{\frac{\pi \sigma}{2}}
		\left(\pi \cos(\sigma v) - 2 v \sin(\sigma v)\right)
	}{(e^{\pi \sigma} + 1)(4v^2 + \pi^2)}, \\
	I_M^{\{\tanh\}}(v) &\simeq
	\frac{
		4 e^{\frac{\pi \tau}{2}}
		\left(\pi \sin(\tau v) + 2 v \cos(\tau v)\right)
	}{(e^{\pi \tau} - 1)(4v^2 + \pi^2)}, \\
	I_M^{\{\relu\}}(v) &= \frac{|v|}{2} - \frac{v \mathrm{Si}(\sigma v)}{\pi}
    - \frac{\cos(\sigma v)}{\sigma},
\end{align}
likely have Fourier transforms or approximate Fourier transforms which could be used to reduce them to polynomials in $v$ valid on $(-V, V)$.  Adding these coefficients to the expansion polynomials would
further reduce the error within $(-V, V)$ at the cost of more rapid divergence outside of $(-V, V)$.  Since
$(-V, V)$ can be chosen conservatively for a given application this trade off is acceptable.

Finally, the method for reducing the expanded forms to equivalent representations in terms of more general special functions requires the most new research.  One possible direction would involve expanding special functions such as the Meijer-G itself using the method.  This would allow for a procedure where a given instance of the Meijer-G could be expanded and the expansion results compared to the expansion of a known function such as $\tanh(\tanh(v))$.  Agreement would indicate that the two functions are the same on $(-V, V)$.  If agreement could be shown for all $V$ then an equivalent representation would be proved.  This would at least allow searching for equivalent representations of functions via a systematic guess and check style approach which could be automated, and would not require directly proving the Meijer-G representation or approximate representation (if one indeed exists) of a function such as $\tanh(\tanh(v))$ which could likely require a great deal of \textit{ad hoc} manipulation.  A more ambitious approach still could involve using the expansions as a means to extract representations from the functions to be expanded themselves which reveal the similarity to the structure of the Meijer-G or related function definitions.  Although this is an ambitious goal, the results of the nonlinear amplifier case studies of Chapter \ref{chapter_additional_case_studies} hints at the possibility of pursuing this path.  We note that applying the AMITE technique to the amplifier case study produces forms for which similarity to the Meijer-G structure is apparent, i.e.,
\begin{align}
	\begin{split}
	&\mu_{n', m} =\nonumber\\
	&\ \ \lim_{n \to n'}
	\frac{i \pi  2^{-2 m-3} e^{-i \pi  n} \left(1+e^{i \pi  n}\right) \Gamma (2 m+2) \left(\sec \left(\frac{1}{2} \pi  (2 m+n)\right)-e^{i \pi  n} \sec \left(\pi  m-\frac{\pi  n}{2}\right)\right)}{\Gamma \left(m-\frac{n}{2}+\frac{3}{2}\right) \Gamma \left(m+\frac{n}{2}+\frac{3}{2}\right)}, \\
	&\quad\quad\forall\quad n' \in \mathbb{W},\ m \in \mathbb{W}.
	\end{split} \\
    \begin{split}
	q_m(g,d)
	= \lim_{g \to g'}
	-\frac{\sqrt{\pi } \sqrt{\frac{g}{g-d}} \sqrt{\frac{g}{d+g}} \Gamma (2 m+3) \left((d-g)^{2 m+2} h_1(t) - (d+g)^{2 m+2}h_2(t)\right)}{g (m+1) \Gamma \left(2 m+\frac{5}{2}\right)}.
	\end{split}
\end{align}
Here the intermediate results $h_1(t)$ and $h_2(t)$ are defined in Chapter \ref{chapter_additional_case_studies} in the case study on biased nonlinear amplifiers.  The similarity to the Meijer-G structure is apparent in the ratios of Gamma functions multiplied by powers of intermediate functions.

\section{Future Studies}
Assuming that the challenges of nested nonlinearity can be addressed, the following studies would comprise a small sampling of the work that would be enabled as a result.

\subsection{Neural Systems Verification}
As a direction for future work we first propose further experimentation with the use of various model replication strategies as a form of verification for deeper neural networks, building off the fundamental concept that introducing a replication step into the verification procedure provides an inherent guarantee of test signal coverage, since test signals with insufficient coverage will cause the replication step to fail.  Specifically, this could allow a rigorous method of conformance testing, as defined in Chapter \ref{chapter_definitions}, for deep neural networks.  Furthermore, defenses against replication attacks for neural systems deployed on neuromorphic hardware are presented by Yang et al. \cite{yang_thwarting_2019} and a software based defense against replication attacks for public prediction interfaces is presented by Juuti et al. \cite{juuti_prada_2019}.  It is proposed that with further work replication approaches for verification could be combined with these approaches and with principles from cyber-security to enable replication for the purposes of system verification for authorized users while blocking replication for unauthorized users.

As noted in Chapter \ref{chapter_introduction}, the degree of deviation between actual weights and specified weights which may be considered a defect is dependent on the structure and other weight values of the network being tested.  In some cases, a small deviation in weight values may cause a large deviation in the resulting output values for some inputs, leading to a defective network.  However, in other cases even a large change in weight values may have a numerically insignificant effect on the input-output manifold, again due to intrinsic nonlinearity.  One example is a saturated hyperbolic tangent neuron.  If the neuron is already well past saturation, increases in the weight going into the neuron will have a numerically insignificant effect on the output.  For this reason, a characterization of the accuracy of the conformance test procedure of Chapter \ref{chapter_network_replication} with respect to the degree of deviation from original values for weights chosen at different points in the architecture of different networks would be a useful follow-on result to extract.  Given that different deviations in weight values will have different effects on the input-output manifold, and some will yield approximately equivalent configurations as discussed in Chapter \ref{chapter_network_replication}, this study would have to be conducted systematically over many fundamental (e.g., very small) and practical (e.g., wide and deep architectures) networks to extract a full characterization of the sensitivity of the test.

Furthermore, fuzz testing has also been presented as a means to address neural network verification challenges, particularly for networks implemented in software.  White-box fuzz testing approaches with test coverage maximization strategies are available and demonstrated in recent work \cite{pei_deepxplore_2017,guo_dlfuzz_2018,odena_tensorfuzz_2019,liang_deepfuzzer_2019}.  These test strategies assume access to code and could be applied during software test stages to complement the methods presented here.  As an extension to this paper and to the body of neural network test literature, we propose experimentation into the effectiveness of different test methods when employed at various stages of the system integration process of Fig. \ref{fig_dev_process}, both as standalone tests and in connection with other test methods.  We also propose investigation into the minimum stimulus signals required to guarantee successful model replication for networks of varying sizes.

The experiments of Chapter \ref{chapter_network_replication} reveal in Fig. \ref{fig_accuracy_wrt_entropy} that the accuracy of the replication procedure, equivalent to training a new neural network from the input-output pairs collected from the previous network, is dependent on the entropy of the interrogation signal.  The accuracy steadily increases with respect to the entropy of the input signals until it reaches a limiting value such that further increases in entropy come with diminishing improvements in accuracy.  It is hypothesized that this limiting value would be larger for larger networks.  In a follow on study, the minimum entropy required of a training database to allow direct reliable replication of a network (i.e., to allow training a new network with an equivalent AMITE polynomial representation) could be determined empirically by regenerating Fig. \ref{fig_accuracy_wrt_entropy} for increasingly large networks.

The point on the accuracy vs. entropy curve at which the improvements in accuracy diminish (e.g., accuracy reaches 0.99 such that further improvements are of diminishing significance) could then be noted for each increase in network size such that a new curve representing minimum entropy required to achieve the specified accuracy could be exposed.  The accuracy threshold could then be varied such that a surface showing minimum entropy required to achieve each accuracy benchmark for each network size could be uncovered.  It is hypothesized that each bisection of this surface created by fixing the accuracy parameter (thereby specifying a curve representing the minimum entropy required to train a network of the given size to the given accuracy) would be monotonically increasing and thus invertible.  The inverse curve would therefore provide insights into the minimum network size required to fit a function to a data set having the given entropy.

It is further hypothesized that both the information content of the input to the function to be learned and the information content of the corresponding output would affect the complexity of the network required to fit the data.  It is therefore further proposed to empirically analyze surfaces showing the network size required to fit functions of varying input and output entropy, and to generate a family of these surfaces for different accuracy benchmarks; e.g., each surface represents the network size required to fit a function having a given input entropy and output entropy to the specified accuracy.  These surfaces would comprise the results of another analysis, distinct from the first analysis and revelatory of different information.  The first analysis determines the minimum input entropy required to allow repeatable replication of the target network as an AMITE polynomial, and as such, is designed to extract knowledge about the minimum entropy required to allow replicable fitting for a given network architecture, as divorced from the characteristics of the function to be learned.  In contrast, the second analysis is designed to allow extraction of knowledge about the minimum network size required to fit datasets having given input and output entropy.

Finally, it is conjectured that the minimum network size is likely a monotonically increasing function of the information gained by concatenating each output value to its corresponding input value; i.e., the minimum network size is a monotonically increasing function of the information gain as defined by the difference in the entropy of the dataset having concatenated inputs and outputs such that if all outputs equal all inputs, the increase in entropy is equal to zero.  In such a case it is expected that the required network is an identity function, i.e., a network of size one.  Any addition of entropy to the concatenated input-output data set over the input data set will demand a more complex model and as a result demand a larger network.

Once these empirical analyses have been conducted, it is proposed that AMITE polynomials could be used as a means to attempt to theoretically derive the curves for (1) entropy required to allow replication of a network of a given size with given accuracy (i.e., entropy of replication stimulus with respect to network size) and (2) network size required to fit a function of a given information gain from input to output as defined above.  If successful, and if agreement between the empirical and theoretical curves could be confirmed, such a study would have significant impacts on the growing field of information theoretic learning.

\subsection{Neural Network Interpretability and Explainability}
If a method for turning deeply nested nonlinear functions into analyzable forms can be derived, a number of experiments in the explainability and interpretability of neural networks could be conducted.  Due to the fundamental limit of the universal approximation theorem, i.e., that the network is only guaranteed to approximate the function to be fitted on a bounded interval, there is never any guarantee that a network's learned function will generalize outside of the feature space or volume wherein training data was provided.  This problem is especially important when using a neural network to learn to solve problems for which the underlying physics leaves a relationship to be found, albeit a relationship which may not be simply expressed in known mathematical language.  The network will fit the function in the space where it was trained but will have no guarantee of fitting the function outside that region.

Given that the nested nonlinear form eludes analysis, it is difficult to extract from the network the function or algorithm that was learned to prove it will not exhibit an undesired result outside this region.  Furthermore, if the network did learn a simple relationship expressible in terms of elementary functions, there are no means of proving this and no means of determining if it is true for a wide variety of cases.  It is proposed that analytically turning the nested nonlinear functions which define the network's inference processing into analyzable forms such as a polynomial via direct expansion using an extended AMITE method, or expansions which can be converted to a form such as Meijer-G, could provide a means to solve this problem.  A simple experiment where a network is trained to fit a simple function and then the function is recovered from the network's weights via analysis would demonstrate that this could be done.  Once demonstrated, such a method could be used to arbitrate between multiple networks, trained using a population based approach, to identify which has the best chance of generalizing outside the training region.  This could also help arbitrate between networks which learn relationships that meet other constraints, such as physics constraints or safety requirements.

\subsection{Scalability for Analysis of Larger Networks}

While some scalability of the techniques developed herein has been demonstrated thus far, it is proposed that further enhancements could increase the scope of the neural network architectures to which these techniques would be applicable.  In some applications, such as neural network equivalency testing as demonstrated in Chapter \ref{chapter_network_replication}, it is necessary to convert an entire network to a polynomial.  However, in many other applications, such as neural network range bounding, this is not required.  A distinction must therefore be made between future studies wherein an entire network must be converted to a polynomial form and future studies where only the individual activation functions or individual layers must be converted to a polynomial form and the handling of nested forms is not required.  In the former case where a single polynomial form is required, it is likely that techniques for directly converting nested forms to practically sized polynomial representations will need to be developed.  For example, to convert a deep neural network to a single polynomial in a fully analytic manner, the methods described in this outline of future work will need to be implemented directly.

However, in many cases such methods will not be required.  In considering the second example, neural network range bounding, only each activation function individually must be converted to a polynomial.  The propagation of errors through each nonlinear function, once represented as a polynomial with formal error bounds, can be handled via the well-known methods of Taylor modeling.  While Taylor models are usually constructed from Taylor series, other polynomials such as Chebyshev, and now the AMITE polynomials with the developments presented herein, can also be used where their advantageous properties are required.  A future study could therefore be conducted where each nonlinear function in a neural network could be converted to a Taylor model via AMITE polynomials and their formal error bounds such that the Taylor models could be composed to efficiently extract a Taylor model of a deep neural network.

It is hypothesized that an AMITE polynomial constructed Taylor model would have smaller error, smaller range over-estimation, and enable modeling across wider intervals than a standard Taylor model.  Since the methods presented herein require only that the nonlinear functions to be converted to AMITE polynomials have generalized Fourier transforms, all of the common operations employed in state of the art neural networks, including convolutions, ReLUs, and even custom activation functions could be treated as per the methods developed herein.  In doing this a catalog of AMITE polynomial representations of many of the nonlinear forms arising in neural networks could be developed and added to the existing repertoire of polynomial expansions as standard reference material to facilitate this analysis.

While a number of computational optimizations, including GPU parallelized matrix operations and custom function lookup tables with custom optimized memory management schemes, have been implemented for the AMITE polynomials and the neural network analysis methods thereby enabled (described in Chapter \ref{chapter_network_replication}), is is proposed that a number of further optimizations could be made, especially for use cases where only individual nonlinear functions must be expanded as polynomials.  First, the computation of the AMITE polynomials for each nonlinear function could be managed via a threading scheme, such that both the benefits of sub-processes and GPU parallelization could be exploited simultaneously to perform the initial polynomial conversion.  Second, formal intervals of uncertainty and functions for propagating these through a neural network could be implemented as types compatible with popular deep learning frameworks such as PyTorch.

This would allow both (1) optimization over intervals of uncertainty (which would prove useful when automating the construction of Taylor models from the AMITE polynomials) and (2) propagation of the intervals of uncertainty through the network to effectively range bound the output of the network through only use of standard PyTorch operations.  Finally, it is further proposed that the exploitation of the closed form error approximations for the AMITE polynomials alongside the exact error formulas presented herein would uniquely enable efficient construction of Taylor models with tighter error bounds than possible with Taylor series.  This could be done by using the closed form approximations for the error formulas to first perform a course search for the maximum value of the error on the interval of interest, before performing a fine search using the exact formulas, informed by the known frequency of the error oscillation and locations of error extrema. 

\subsection{Waves and Radar Clutter Analysis}

In the shallow water wave study of Chapter \ref{chapter_additional_case_studies}, the main motivation is to provide a means of analyzing the shape and height of ocean waves traveling in shallow water so their effects on radar signals can be analyzed.  The problem is formulated such that the height of waves with a wavenumber uniformly distributed in a certain interval can be determined at a fixed point in time.  In the context of this work, this problem was chosen to be an exemplary case study to illustrate utility of the AMITE technique.  However, these results could be employed in a much more detailed and complete study on the effects of shallow water waves on naval radar signals.

In such a study, a model of shallow water wave height as a function of space and time could be developed given measured statistics about the waves, e.g., distribution of wavenumers measured over time.  Then, given information about the radar characteristics, analysis could be performed to determine the percentage of time the wave crests intersect the Fresnel zones of the link between the radar and true targets, allowing the extraction of statistics about the percentage of time during which multi-path effects could cause interference.  In a more detailed analysis, the location of true targets and the statistics of wave height could be used to further derive statistics for how often a true target is expected to be obscured and for how often a false alarm due to wave cresting is expected to be observed.  Finally, in an even more detailed study, it is proposed that the effects of the wave shape on the reflection of the radar's pulses could be analyzed to design better filters for removing false alarms.

\subsection{Nonlinear Circuit Design and Analysis}

Finally, a simple yet impactful study is proposed for future work employing the methods of this thesis to the analysis of nonlinear circuits.  In Chapter \ref{chapter_additional_case_studies}, two nonlinear circuits are analyzed in detail to reveal exact relationships governing the properties of their output as a function of input parameters, i.e., gain and bias.  This study could be extended in two main future studies.  First, cases where multiple stages of nonlinear amplification are employed could be considered.  In this study, exact expressions for strength of harmonics at the output of a multistage nonlinear amplifier could be derived as a function of input gain and gain at intermediate stages.  This study could be extended by considering input bias and intermediate stage bias.  These results would be significant for characterizing nonlinear amplifiers, e.g., for deriving expressions by which different instantiations of the same nonlinear amplifier design could be distinguished by analyzing signal artifacts due to slight variations in bias point and gain staging.

A second study in this domain would involved deriving Gibbs-phenomenon reduced expansions of the nonlinear functions in the amplifiers' gain stages.  In Chapter \ref{chapter_improved_expansions}, the exact expressions for the expansion errors are obtained, and from these exact expressions closed form expressions for the Gibbs-like errors in the expansions are derived.  It is reasonably hypothesized that these expressions are of suitable form such that the AMITE technique could be reapplied to them; i.e., the AMITE technique could be applied recursively such that the errors within the convergent region are characterized by a polynomial which is then absorbed back into the coefficients of the original polynomial.  This would enable a precise characterization of the nonlinear effects of a multistage amplifier without propagation of errors through each stage of nonlinearity.  The results could be used for applications where it is necessary to characterize precisely the distortion introduced by a nonlinear amplifier in order to distinguish between multiple nonlinearly distorted signals or to mitigate the effects of nonlinear distortion.  Such analysis would have applications wherever multistage nonlinear amplifiers are employed, i.e., in the analysis of radar signal processing, music and audio signal processing, and optical communication channels to name a few.

\subsection{Fundamental Mathematical Studies in Function Approximation}

As a fundamental study, it is proposed to perform further investigation into the ability of neural networks combined with AMITE polynomials to act as a powerful function approximation method.  This study is motivated by the results of Chapter \ref{chapter_network_replication}, specifically Fig. \ref{fig_accuracy_wrt_snr2} which shows that with a judiciously chosen error metrics a highly accurate test can be achieved even with very low signal to noise ratios in the measurements of the outputs which occur in response to the interrogation signal.  Noting that before the successful algorithms of Chapter \ref{network_replication} were developed, attempts were made to fit polynomials directly to the input-output pairs, and further noting that this failed due to ill conditioning of the matrix which must be pseudo-inverted to do so hints at why the successful method performed so well against noise.  While pseudo-inverting a matrix requires optimization over a large amount of free parameters, fitting a network directly and then converting the network to a polynomial form analytically through use of the AMITE polynomials allows the nested structure of the network to provide sufficient fitting power to capture the details of the function defined by the input-output pairs without needing as many free parameters as the polynomial directly.

Since converting from a neural network to an AMITE polynomial is a lossless and fully analytic process, the benefits of fitting a structure with fewer free parameters (but less analytic convenience than a polynomial) and then converting to a convenient polynomial form are fully realized.  This result is expected, since it is now well understood in the mathematics and artificial intelligence engineering communities that nested models, such as neural networks, fit functions more efficiently than flat models since the early layers learn useful intermediate features.  To exploit this phenomenon fully, further study is proposed on the extraction of function approximations by first fitting a neural network or other nested nonlinear form to data and then extracting a polynomial analytically from the result.  It is further conjectured that the free parameters of such models would be useful features themselves in the analysis of dynamical systems which could then be utilized as inputs in neural or conventional decision making processes and then analyzed with rigor via the AMITE polynomials.